% !TeX root = ../0_Manuscript.tex

\section{Body Biasing Injection enhanced practice \ddc}
\label{chap2:goodPractices}
With the platform thoroughly introduced,we can now discuss the BBI practice in the state of the art compared to the enhancements we propose.
\subsection{BBI practice in the state of the art \ddc}
% !TeX root = ../0_Manuscript.tex

%\begin{figure}[H]
%    \centering
%    \includegraphics[width=0.7\textwidth, center]{2_goodPractices/figures/model0Probe.pdf}
%    \caption{Simplified model of a \bbi platform in the state-of-the-art. The model represents the major tools involved in a \bbi platform, with the voltage pulse generator, the \bbi metal probe, the transmission line used to interconnect the generator to the probe, the IC, modeled with a 1 k\textOmega\xspace resistor, and the platform grounding impedance, represented with two 150 \textOmega\xspace resistors. This model allows for fast evaluation and prediction of the \bbi platform macro-behavior.}
%    \label{fig:bbiPracticeBadGnd}
%\end{figure}

\begin{figure}[ht]
    \centering
    \begin{subfigure}{0.45\textwidth}
        \includegraphics[width=1.0\textwidth, center]{2_goodPractices/figures/model0Probe.pdf}
        \caption{Model of a state-of-the-art \bbi setup, with its key components.}
        \label{fig:bbiPracticeBadGnd}
    \end{subfigure}
    \hspace{0.045\textwidth}
    \begin{subfigure}{0.45\textwidth}
        \includegraphics[width=1.0\textwidth, center]{2_goodPractices/figures/SIMPLE_BBI_0.pdf}
        \caption{Simulation results of the state-of-the-art \bbi platform model.}
        \label{fig:bbiPracticeBadGndSignals}
    \end{subfigure}
    \caption{State-of-the-art \bbi platform electrical model and its simulation results.}
    \label{fig:bbiBadGndGlobalFig}
\end{figure}

First and foremost, in order to introduce the proposed enhancements, we will explain the limitations of state of the art BBI practices.
Fig. \ref{fig:bbiPracticeBadGnd} presents an example BBI platform as it is performed in the state of the art.
This schematic only represents the main pieces of equipment, roughly composed of:
\begin{itemize}
    \item A voltage pulse generator
    \item A transmission line
    \item A BBI probe
    \item An IC target
    \item A grounding setup
\end{itemize}
In that case, we want to bring forward two limitations of the platform as it stands.
The first one is that most commercially available high voltage generator require a very specific load impedance to be used in order to meet their specifications, thus the \mbox{50 \textOmega} transmission line.
However, the vast majority of integrated circuits substrates do not present at all times and in every location the exact required impedance for a specific generator.
% !TeX root = ../0_Manuscript.tex

\begin{figure}[H]
    \centering
    \includegraphics[width=15cm, center]{2_goodPractices/figures/SIMPLE_BBI_0.pdf}
    \caption{BBI platform example in the state of the art: simulation}
    \label{fig:bbiPracticeBadGndSignals}
\end{figure}
Therefore, impedance mismatch effects will inevitably show up, such as bounces in the transmission line.
These effects are illustrated in Fig. \ref{fig:bbiPracticeBadGndSignals}.
As one can notice, due to the bounces, there is a 25 \% undershoot on the pulse received by the target (green signal), in addition to a 100 V overshoot.
It is directly echoed on the IC ground current measurement (purple signal).
However, one might wonder why this can be a problem for the practice of BBI.
The bounces in the transmission line, reflected in the IC current, are then propagated through the IC.
Therefore, because the IC does not present in every substrate location a constant impedance, it is impossible to predict the bounces' shape.
Thus, because an attacker has to achieve precise timing when injecting faults into an IC, and because of this impossibility to reach a certain degree of precision, it is highly problematic in a fault injection context.
In order to mitigate this problem without increasing the platform's cost, we propose a very simple addition to the platform, consisting in connecting at the voltage pulse generator output a \mbox{50 \textOmega} resistor in order to place the generator closer to its specifications.

The second limitation concerns the platform grounding.
%The example platform of Fig. \ref{fig:bbiPracticeBadGnd}
Contrary to impedance mismatching, it is a less predictable issue, as each platform will have very different grounding installation.
Therefore, it is important to take measures to limit its side effects.
In addition to worsening impedance mismatching, it can lead to a non-negligible limitation concerning energy transfer to the IC, by acting like a resistive voltage divider.
Alleviating this limitation is simpler than matching the load impedance, as it only requires low-length, low-resistance ground interconnections to be made on top of the existing ones.

\subsection{Enhanced BBI practice \ddc}
% !TeX root = ../0_Manuscript.tex

\begin{figure}[H]
    \centering
    \includegraphics[width=17cm, center]{2_goodPractices/figures/model2ProbeManuscrit.pdf}
    \caption{The proposed enhanced BBI platform: schematic}
    \label{fig:bbiBestPractice}
\end{figure}
% !TeX root = ../0_Manuscript.tex

\begin{figure}[H]
    \centering
    \includegraphics[width=\textwidth, center]{2_goodPractices/figures/SIMPLE_BBI_2.pdf}
    \caption{The proposed enhanced BBI platform: simulation}
    \label{fig:bbiBestPracticeSignals}
\end{figure}
Fig. \ref{fig:bbiBestPractice} and Fig .\ref{fig:bbiBestPracticeSignals} both illustrate the proposed improvements.
The red nets on the schematic indicate the enhancements, and on the simulation results are shown the observed signals in the state of the art on top of the new improved signals.
What is important to remark here is first the significant bounces' reduction both on the IC received voltage pulse waveform and on the IC ground current waveform (which is an image of the received voltage due to the purely resistive nature of the simulated IC).
Then, the effective injected current stays the same between both scenarios, which is expected.
First because reducing the grounding impedance allows an increase in current, second because the approximate impedance matching reduces the injected current over time due to bounces' reduction.

\subsection{BBI practices: a simple actual comparisons \ddc}
Evaluating the effects of the proposed enhancements through simulations is the first step.
In order to verify their soundness, we will analyze in the first place the same scenario as before with real components.

%The first noticeable one is the imperfect electrical grounding.
%Indeed, on the example platform, one can note the 150 \textOmega ground impedance between each piece of equipment ground and the reference ground.
